% BHU PYQ -- Auto-generated & error-corrected
\documentclass[11pt,a4paper]{article}

\usepackage[a4paper,top=2.5cm,bottom=2.5cm,left=2.8cm,right=2.8cm,headheight=14pt]{geometry}
\usepackage{amsmath,amssymb}
\usepackage[T1]{fontenc}
\usepackage[utf8]{inputenc}
\usepackage{lmodern}
\usepackage{microtype}
\usepackage{fancyhdr}
\usepackage{enumitem}
\usepackage{hyperref}
\usepackage{lastpage}
\usepackage{parskip}

\hypersetup{colorlinks=true,urlcolor=black,linkcolor=black,
  pdftitle={STA/STB-501: Applied Statistics -- BHU PYQ},pdfauthor={Banaras Hindu University}}

\pagestyle{fancy}\fancyhf{}
\renewcommand{\headrulewidth}{0.4pt}\renewcommand{\footrulewidth}{0.4pt}
\fancyhead[L]{\small\textit{Banaras Hindu University}}
\fancyhead[R]{\small\textit{STA/STB-501: Applied Statistics}}
\fancyfoot[C]{\small Page \thepage\ of \pageref{LastPage}}

\newlist{parts}{enumerate}{2}
\setlist[parts,1]{label=(\alph*),leftmargin=*,itemsep=5pt,topsep=3pt}
\setlist[parts,2]{label=(\roman*),leftmargin=*,itemsep=3pt,topsep=2pt}
\newcommand{\pts}[1]{\hfill\textit{[#1]}}

\begin{document}

\begin{center}
  {\large\bfseries BANARAS HINDU UNIVERSITY}\\[2pt]
  {\normalsize B.A./B.Sc. (Hons.) Semester V Examination, 2016-17}\\[6pt]
  \rule{0.8\linewidth}{0.5pt}\\[6pt]
  {\large\bfseries Paper No. STA/STB-501: Applied Statistics}\\[4pt]
  {\normalsize Time: 3 Hours\hfill Maximum Marks: 70}
\end{center}
\rule{\linewidth}{0.4pt}
\smallskip
\noindent\textit{Instructions: Answer any \textbf{five} questions. All questions carry equal marks. Symbols have their usual meanings.}
\bigskip

\medskip\noindent\textbf{Question 1.}
\begin{parts}
  \item Discuss GRR and NRR in detail with suitable examples. State the interpretation of the statement ``NRR $= 1$''. \pts{7}
  \item Describe the general fertility rate and total fertility rate with examples. Write the interpretation of TFR $= 2$. \pts{7}
\end{parts}

\medskip\rule{\linewidth}{0.2pt}

\medskip\noindent\textbf{Question 2.}
\begin{parts}
  \item What are the methods of obtaining Vital Statistics? What are the various sources of Vital Statistics? Also give the uses of Vital Statistics. \pts{7}
  \item Give a brief history of the Census of India. Who was the Registrar General for Census 2011? \pts{7}
\end{parts}

\medskip\rule{\linewidth}{0.2pt}

\medskip\noindent\textbf{Question 3.}
\begin{parts}
  \item What do you mean by a life table? Write the assumptions of the life table. Also discuss the various columns of a complete life table. \pts{7}
  \item Discuss all the important curves used for measuring the growth of a population (e.g., exponential, logistic). \pts{7}
\end{parts}

\medskip\rule{\linewidth}{0.2pt}

\medskip\noindent\textbf{Question 4.}
\begin{parts}
  \item Discuss the importance and limitations of index numbers. \pts{7}
  \item Define crude death rate. How does standardisation help in overcoming the limitation of the crude death rate? Explain clearly the standardised death rates, including the methods of their computation. \pts{7}
\end{parts}

\medskip\rule{\linewidth}{0.2pt}

\medskip\noindent\textbf{Question 5.}
\begin{parts}
  \item What is a price index number? Explain the meaning of the current year and base year in reference to an index number. \pts{7}
  \item What is meant by a cost of living index number? Describe the procedure followed in its preparation. \pts{7}
\end{parts}

\medskip\rule{\linewidth}{0.2pt}

\medskip\noindent\textbf{Question 6.}
\begin{parts}
  \item Discuss the various tests of index numbers and verify that Fisher's index number is an ideal index number. \pts{7}
  \item What do you mean by time series? What are its uses? \pts{7}
\end{parts}

\medskip\rule{\linewidth}{0.2pt}

\medskip\noindent\textbf{Question 7.}
\begin{parts}
  \item Write short notes on the different types of components of time series. Explain the difference between seasonal and cyclical variations with an example. \pts{7}
  \item Explain secular trend in detail. Critically examine the various methods of measuring trend. \pts{7}
\end{parts}

\medskip\rule{\linewidth}{0.2pt}

\medskip\noindent\textbf{Question 8.}
\begin{parts}
  \item What is statistical quality control (SQC)? Discuss the need and utility of SQC in industry. Explain why $3\sigma$ limits are used in SQC charts. \pts{7}
  \item Discuss in detail the construction of the $p$-chart and $c$-chart. What is the theoretical justification for using the Poisson distribution in the case of the $c$-chart? \pts{7}
\end{parts}

\medskip\rule{\linewidth}{0.2pt}

\medskip\noindent\textbf{Question 9.}

Discuss the construction of the $(\bar{X}, R)$-chart, explaining the role of control limits and the rationale for using the range $R$ as a measure of variability. \pts{14}

\vfill
\rule{\linewidth}{0.4pt}
\begin{center}\small
  \href{https://bhu-exam.vercel.app/}{Click here to visit our website}\\[4pt]
  \href{https://me-aryan.vercel.app/}{Click here to visit our contributor}\\[4pt]
  \href{https://quashan-me.vercel.app/}{Click here to visit our contributor}
\end{center}

\end{document}